\chapter{\MakeUppercase{Project Description and Goals}}
\justifying

\section{Literature Review}

A systematic review of 49 research papers was conducted across three disease domains — cardiac arrhythmia, diabetes mellitus, and Parkinson's disease — to identify the state of the art, assess existing approaches, and identify the research gaps that motivate this project. The papers are organised by domain below.

% ──────────────────────────────────────────────────────────────────────
\subsection{Cardiac Arrhythmia Detection using ECG}
% ──────────────────────────────────────────────────────────────────────

\textbf{Review of existing methods.}
Ayyub et al. \cite{ayyub2025} conducted a systematic review of 219 papers on AI-based arrhythmia detection, noting that reported accuracies of individual models range from 50\% to 99.9\%, but the models are overwhelmingly evaluated on a single dataset (MIT-BIH) and suffer from severe class imbalance, particularly for rare beat classes (S, F).

Several papers have proposed one-dimensional \ac{cnn} architectures as the dominant approach. Guhdar et al. \cite{guhdar2024} proposed a 1D-CNN with channel and spatial attention achieving 99.3\% accuracy on MIT-BIH (5-class AAMI). Kavitha and Shakkeera \cite{kavitha2024} achieved 99.5\% with an explainable 1D-CNN using LIME/SHAP, but noted that the interpretability methods are computationally expensive and have not been validated by clinicians. Kerdoudi et al. \cite{kerdoudi2024} achieved 99.5\% with a multi-head attention CNN-BiLSTM architecture, and Venkatesh et al. \cite{venkatesh2026} proposed ASNet (Attention-Stacked Network) achieving 99.2\%.

Multi-scale and multi-modal architectures have also been explored. Khatar et al. \cite{khatar2024} proposed GAF-GradCAM, which encodes ECG signals as Gramian Angular Field images and performs dynamic weighted fusion, achieving 99.4\% accuracy, but at the cost of increased memory and computational requirements. Mahajan et al. \cite{mahajan2024} integrated CNN, GNN, and BiLSTM for multi-lead ECG achieving 99.1\%, but requiring multi-lead input which is less accessible. Tang et al. \cite{tang2024} designed LEPENN, a lightweight model with only $\sim$15K parameters achieving 98.9\%, specifically targeting embedded deployment. Moaven and Abbaszadeh \cite{moaven2026} proposed a Hybrid PCA-PINN framework achieving 98.7\%, requiring domain-specific physical constraint design.

Transformer-based approaches have emerged as a second major direction. Dutenhefner et al. \cite{dutenhefner2026} proposed a Hybrid Hierarchical Transformer combining CNN and Transformer encoders achieving $\sim$97\% on MIT-BIH but requiring 12-lead ECG data. Liu et al. \cite{liu2024} developed a Transformer-Convolutional network with Grad-CAM achieving 98.2\% on MIT-BIH but with high computational requirements. Duranay et al. \cite{duranay2024} applied a Vision Transformer to ECG images achieving 97.3\%, but highlighted that image conversion loses temporal resolution.

For patient-specific classification, Prakash and Atef \cite{prakash2024} demonstrated 99.3\% accuracy with a lightweight CNN-LSTM capturing local and long-term dependencies, but noted that patient-specific models require initial labelled data for each new patient.

Sleep-apnea-related ECG analysis was explored by Hou et al. \cite{hou2024sleep} and Hou et al. \cite{hou2024tho}, both achieving $\sim$95\% on the Apnea-ECG database using CNN-Transformer-LSTM hybrids, though these focus on sleep apnea rather than arrhythmia classification.

\textbf{Key limitations identified across ECG literature:}
\begin{itemize}
  \item All high-performing models (99\%+) are trained and tested on the same benchmark (MIT-BIH), with no cross-dataset evaluation.
  \item Rare \ac{aami} classes — Supraventricular (S) and Fusion (F) — consistently achieve F1-scores below 0.20 due to extreme class imbalance ($<$2\% of beats).
  \item Transformer and multi-lead models are computationally infeasible for CPU deployment.
  \item No existing system integrates ECG arrhythmia detection within a conversational health screening platform.
  \item High accuracy on MIT-BIH does not guarantee clinical utility due to train-test data leakage within the same patient cohort.
\end{itemize}

% ──────────────────────────────────────────────────────────────────────
\subsection{Diabetes Mellitus Prediction}
% ──────────────────────────────────────────────────────────────────────

\textbf{Review of existing methods.}
Diabetes prediction from clinical biomarkers has been extensively studied using the Pima Indians Diabetes Dataset (PIDD). Classical ensemble and optimisation-based approaches dominate the literature.

Daza et al. \cite{daza2024} used a multi-level stacking ensemble (RF, SVM, NB, KNN, LR as level-1; Gradient Boosting as meta-learner) achieving 82.8\% on PIDD and 85.3\% on the Frankfurt Hospital dataset. Reza et al. \cite{reza2024} used a similar stacking ensemble achieving 81.2\% on PIDD. Gollapalli et al. \cite{gollapalli2022} proposed a stacking ensemble on a Saudi Arabian 3-class dataset achieving 94.48\%, but on a population-specific dataset not publicly available. Elgendy et al. \cite{elgendy2025} proposed a dual-stage explainable ensemble using SHAP feature selection and stacking (RF, XGBoost, LightGBM, SVM) achieving 82.5\% on PIDD and 86.1\% on BRFSS. Phan et al. \cite{phan2025} proposed REMED-T2D with Bayesian hyperparameter optimisation, achieving 83.9\% on PIDD, 98.5\% on the Early Stage Diabetes dataset, and 86.2\% on BRFSS.

Metaheuristic feature selection methods have been paired with classical classifiers. Siva Shankar and Manikandan \cite{sivashankar2019} used Grey Wolf Optimization (GWO) with Fuzzy Logic achieving 83.12\%. Vekkot et al. \cite{vekkot2025} extended GWO with SHAP explainability achieving 82.5\%. Sirmayanti et al. \cite{sirmayanti2024} proposed an autophagy-inspired GWO variant achieving 83.7\%. Karuppasamy et al. \cite{karuppasamy2025} combined PSO and Genetic Algorithms (PG-S approach) achieving 84.5\%. Naskar et al. \cite{naskar2025} proposed a Guided Population-based Genetic Algorithm achieving 79.2\% on PIDD. Basith et al. \cite{basith2024} used Particle Swarm Optimization for feature fusion achieving 88.3\% on PIDD.

Vasuki et al. \cite{vasuki2024} built a Knowledge-aware Attentional Neural Network optimised with WV-GWO achieving 87.3\% on diabetes prediction. Bhagat and Maulik \cite{bhagat2024} added Blockchain to a stacked ensemble (StaBloCare) achieving 86.4\% but with significant infrastructure overhead. Wang and Nguyen \cite{wang2024} proposed a Wasserstein Autoencoder for synthetic tabular data generation, improving downstream classification by 2\%--8\% through data augmentation. Chen et al. \cite{chen2026} used Enhanced Adaptive Differential Evolution for feature selection achieving 84.2\% with SVM.

\textbf{Key limitations identified across diabetes literature:}
\begin{itemize}
  \item All studies are trained and tested on the PIDD (768 samples, female patients of Pima Indian heritage aged $\geq$21), a small, highly demographic-specific dataset with known bias.
  \item Accuracies plateau below 90\%, with the best published results ranging from 81\%--88\% on PIDD. No paper significantly surpasses this ceiling.
  \item No existing system collects the 8 diabetes biomarkers \textit{conversationally} through natural language dialogue; all prior systems require structured form-based data entry.
  \item Metaheuristic approaches (GWO, GA, PSO) add significant computational overhead with only marginal accuracy improvement over simpler methods.
  \item The diabetic class (positive) consistently suffers from lower recall, meaning many diabetic patients are missed — a critical failure mode for a screening tool.
\end{itemize}

% ──────────────────────────────────────────────────────────────────────
\subsection{Parkinson's Disease Detection}
% ──────────────────────────────────────────────────────────────────────

\textbf{Review of existing methods.}
Voice-based Parkinson's detection builds on the seminal UCI Parkinson's dataset introduced by Little et al. \cite{little2008}, which established 22 acoustic biomarkers (jitter, shimmer, HNR, nonlinear dynamics) as powerful discriminators of PD.

Guatelli et al. \cite{guatelli2024} applied Extreme Learning Machine (ELM) with random weight neural networks to Mel spectrograms of voice recordings, achieving 96.5\% accuracy. Palakayala and Kuppusamy \cite{palakayala2024} combined a Deep Structured Neural Network with a Genetic Algorithm for feature optimisation, achieving 97.6\% binary and 93.7\% multi-class accuracy. Khafaga et al. \cite{khafaga2024} proposed HGGO-PSO (hybrid metaheuristic) with SVM achieving 97.4\% on UCI. Chakraborty et al. \cite{chakraborty2025} achieved 97.4\% with voice alone, and 98.6\% with multi-modal fusion (voice + gait + handwriting).

Beyond voice, alternative modalities have been explored. Majhi et al. \cite{majhi2024} applied transfer learning with InceptionV3 to spiral/wave drawing images achieving 89.75\%. Hammouda et al. \cite{hammouda2024} applied ResNet-based transfer learning to time-series imaging from wrist-worn IMU sensors achieving F1 = 0.986. Wang et al. \cite{wang2024point} developed PointExplainer, a point cloud neural network for 3D motion capture data achieving 92.3\%. Gao et al. \cite{gao2024} applied multi-scale spatio-temporal networks with topological features from EEG achieving 93.8\%. Soylu et al. \cite{soylu2025} developed a CNN-RNN hybrid for multi-disease neurological detection including PD, achieving 96.8\% accuracy.

Speech signal analysis is reviewed by Rusz et al. \cite{rusz2024}, noting that digital speech biomarkers can detect PD at prodromal stages (65\%--85\% accuracy). Peixoto et al. \cite{peixoto2024} used multi-agent LLM collaboration (GPT-4, Med-PaLM) for multimodal PD diagnosis achieving 89.5\%, but relying on proprietary APIs and non-public clinical data. Wagdarikar and Jagtap \cite{wagdarikar2025} proposed a two-way voice feature representation (TWVFR) with 1D+2D DCNN fusion achieving 98.33\% on the Saarbruecken Voice Database.

Review papers provide broader perspective: Islam et al. \cite{islam2024} reviewed ML/DL for PD detection from handwriting and voice (75\%--99\% across studies); Chandru et al. \cite{chandru2024} reviewed AI for neurodegenerative diseases; Zhao et al. \cite{zhao2024} surveyed AI-enabled PD detection via multimodal data; Dixit et al. \cite{dixit2024} provided a comprehensive review across methodologies, datasets, and clinical applications; Sigcha et al. \cite{sigcha2023} systematically reviewed DL + wearable sensors for PD (80\%--99\%); Lee and Lu \cite{lee2025} reviewed clinical interventions (DBS, MRgFUS) and Shi et al. \cite{shi2026} reviewed MOF/COF drug delivery with AI for PD treatment.

\textbf{Key limitations identified across Parkinson's literature:}
\begin{itemize}
  \item The UCI Parkinson's dataset contains only 195 recordings from 31 subjects (23 PD, 8 healthy) — grossly insufficient for statistically robust evaluation.
  \item High accuracy (97\%+) on this dataset is achievable with almost any modern classifier, suggesting subject-level overfitting rather than true generalisation.
  \item Multi-modal approaches (EEG, IMU, point cloud) require specialised, expensive hardware unavailable to most users.
  \item LLM-based approaches \cite{peixoto2024} depend on proprietary APIs, cloud processing of sensitive voice data, and hallucination-prone pipelines with no quantitative model backing.
  \item There is no deployable system that extracts voice biomarkers from a simple WAV file upload, runs them through a DNN, and delivers results conversationally in patient-friendly language.
\end{itemize}

% ──────────────────────────────────────────────────────────────────────
\section{Research Gap}
% ──────────────────────────────────────────────────────────────────────

Based on the systematic review above, the following specific, interconnected research gaps were identified — each of which is directly addressed by this project:

\paragraph{Gap 1: No multi-disease early screening platform exists.}
All 49 reviewed papers focus on a \textit{single} disease domain. No system simultaneously screens for cardiac arrhythmia, diabetes, and Parkinson's disease in a single patient session. This forces patients to seek out multiple isolated tools, reducing the likelihood of comprehensive early-stage risk assessment.
\textbf{$\rightarrow$ Addressed by:} The AI Health Screening Assistant integrates three independent PyTorch models (ECGNet, DiabetesNet, ParkinsonNet) that run simultaneously in a single pipeline, producing a unified multi-domain risk report.

\paragraph{Gap 2: Conversational data collection is completely absent.}
Every reviewed system requires structured form-based data entry. The 8 diabetes biomarkers (glucose, blood pressure, BMI, etc.) must be entered manually by the user in a rigid form. This creates a barrier for non-technical users and reduces data quality.
\textbf{$\rightarrow$ Addressed by:} The Google Gemini \ac{llm} conducts a natural language medical intake interview, collecting all required data through guided conversation (2--3 questions per turn), explaining the purpose of each biomarker in lay language.

\paragraph{Gap 3: LLM hallucination of medical risk values is a patient safety risk.}
Peixoto et al. \cite{peixoto2024} demonstrated LLM-based PD diagnosis, but the system relies on LLMs to generate risk assessments — creating a direct risk of hallucination of clinical values. No reviewed system architecturally prevents the LLM from fabricating probabilities.
\textbf{$\rightarrow$ Addressed by:} The double-pass \ac{llm} architecture strictly separates the LLM's role (data collection and result interpretation) from the quantitative models' role (risk probability generation). The LLM \textit{never} generates risk numbers; all probabilities come exclusively from PyTorch inference.

\paragraph{Gap 4: ECG screening is confined to high-complexity models or multi-lead hardware.}
While high-performing ECG models exist (Transformer, GNN-based), they require multi-lead input \cite{mahajan2024}, significant GPU compute \cite{liu2024}, or proprietary image conversion pipelines \cite{khatar2024}. No efficient single-lead, CPU-deployable arrhythmia classifier exists within a web application.
\textbf{$\rightarrow$ Addressed by:} ECGNet is a lightweight ($\sim$487K parameters) Multi-Scale Attention 1D-CNN that operates on single-lead CSV ECG uploads, runs on CPU, and is deployed within the FastAPI backend.

\paragraph{Gap 5: Voice-based Parkinson's screening is not accessible to ordinary users.}
EEG-based models \cite{gao2024} require clinical EEG equipment. IMU wearable approaches \cite{hammouda2024} require specialised sensors. 3D motion capture \cite{wang2024point} is impractical. Existing voice models evaluate on the UCI dataset but are not deployed in any accessible tool.
\textbf{$\rightarrow$ Addressed by:} ParkinsonNet accepts any WAV recording through a browser upload widget. The librosa-based pipeline extracts all 22 UCI-compatible voice biomarkers in real-time, enabling Parkinson's risk assessment from any microphone recording.

\paragraph{Gap 6: No explainable, patient-facing result delivery in lay language.}
Of the 49 reviewed papers, none provide patient-facing communication of model outputs. Clinician-facing explainability tools (LIME \cite{kavitha2024}, SHAP \cite{elgendy2025}, Grad-CAM \cite{liu2024}) exist, but they produce technical visualisations unsuitable for lay users, and none are integrated into deployable web applications.
\textbf{$\rightarrow$ Addressed by:} After PyTorch inference, the system sends a \texttt{<MODEL\_OUTPUT>} block back to Gemini, which generates an empathetic, non-diagnostic explanation calibrated to the patient's literacy level — with specific action recommendations based on triage level.

% ──────────────────────────────────────────────────────────────────────
\section{Objectives}
% ──────────────────────────────────────────────────────────────────────

The primary objectives of this project, directly motivated by the identified research gaps, are:

\begin{enumerate}
  \item To design and implement a multi-disease \ac{ai} health screening system (Gap 1) covering cardiac arrhythmia, diabetes mellitus, and Parkinson's disease simultaneously.
  \item To develop a conversational medical intake interface (Gap 2) using Google Gemini \ac{llm} that collects patient data through natural dialogue without form-based data entry.
  \item To architect and implement a double-pass \ac{llm} pipeline (Gap 3) that strictly prevents LLM hallucination of medical risk probabilities, with all quantitative values generated by PyTorch models.
  \item To train and deploy ECGNet (Gap 4), a lightweight Multi-Scale Attention 1D-CNN for CPU-deployable single-lead ECG arrhythmia classification.
  \item To implement an end-to-end voice feature extraction pipeline (Gap 5) extracting all 22 UCI-compatible biomarkers from WAV file uploads for Parkinson's risk assessment.
  \item To deliver screening results in patient-friendly, lay language (Gap 6) with automated triage classification and clinically appropriate action recommendations.
\end{enumerate}

% ──────────────────────────────────────────────────────────────────────
\section{Problem Statement}
% ──────────────────────────────────────────────────────────────────────

Despite 49 published research papers demonstrating strong individual machine learning models for cardiac arrhythmia detection, diabetes prediction, and Parkinson's disease assessment, \textit{no accessible, conversational system exists that simultaneously screens a patient across all three domains from a single interaction}. Existing single-disease tools require structured data entry, operate in isolation, produce outputs only interpretable by clinicians, and — in the emerging category of LLM-based tools — risk fabricating clinical risk values through hallucination.

This project addresses this problem by building an end-to-end \ac{ai} health screening assistant that:
(1) collects multi-domain patient data conversationally via a Gemini \ac{llm};
(2) runs three purpose-trained PyTorch neural networks simultaneously for cardiac, metabolic, and motor risk assessment;
(3) prevents LLM fabrication through a strictly structured two-pass pipeline;
and (4) delivers patient-facing, empathetic, non-diagnostic results with automated triage — closing the gap between published ML research and accessible clinical screening tools.

% ──────────────────────────────────────────────────────────────────────
\section{Project Plan}
% ──────────────────────────────────────────────────────────────────────

The project was executed across the academic year 2025--2026 in five phases:

\begin{enumerate}
  \item \textbf{Phase 1 — Literature Survey and Dataset Acquisition} (Weeks 1--4): Review of all 49 research papers across three disease domains; acquisition of MIT-BIH Arrhythmia Database, Pima Indians Diabetes Dataset, and UCI Parkinson's Dataset; gap analysis and system design.
  \item \textbf{Phase 2 — Model Training} (Weeks 5--10): Preprocessing pipelines, ECGNet/DiabetesNet/ParkinsonNet architecture design, training on Google Colab T4 GPU, evaluation and hyperparameter tuning.
  \item \textbf{Phase 3 — Backend Development} (Weeks 11--15): FastAPI orchestration engine, three inference modules (\texttt{heart.py}, \texttt{diabetes.py}, \texttt{parkinson.py}), Gemini \ac{llm} integration with double-pass architecture.
  \item \textbf{Phase 4 — Frontend Development and Integration} (Weeks 16--18): Streamlit \ac{ui}, file upload integration, risk badge display, triage indicator, end-to-end system testing.
  \item \textbf{Phase 5 — Evaluation and Report Writing} (Weeks 19--22): Performance benchmarking against literature, integration testing, error analysis, capstone documentation.
\end{enumerate}

